\title{DeepSeekMath: Pushing the Limits of Mathematical Reasoning in Open Language Models}

\begin{abstract}
Language models often struggle with mathematical reasoning because of the intricate and highly structured demands of this domain. This study presents DeepSeekMath~7B, a model developed by further pre-training DeepSeek-Coder-Base-v1.5 7B on a dataset comprising 120B math-focused tokens collected from Common Crawl, and enriched with both natural language and programming code sources. DeepSeekMath~7B achieves a notable 51.7\% accuracy on the MATH competition benchmark, reaching a level comparable to Gemini-Ultra and GPT-4—accomplished without the aid of external toolkits or voting mechanisms. When using self-consistency over 64 samples, the model attains a 60.9\% score on MATH. The enhancements in DeepSeekMath's mathematical reasoning can be attributed to two principal elements: Firstly, an extensive data selection pipeline designed to fully utilize rich, publicly available web resources. Secondly, the introduction of Group Relative Policy Optimization (GRPO), an adaptation of the Proximal Policy Optimization (PPO) algorithm that both improves mathematical problem-solving capabilities and reduces the memory footprint typically associated with PPO.
\end{abstract}

\section{Introduction}

Large language models (LLM) have transformed methods for mathematical reasoning within artificial intelligence, driving notable progress in tasks measured by the quantitative reasoning benchmark \citep{MATH} as well as the geometry reasoning benchmark \citep{trinh2024solving}. In addition to these achievements, such models have been shown to be valuable aids to humans tackling advanced mathematical challenges \citep{tao}. Despite these successes, the most advanced systems, including GPT-4 \citep{gpt4} and Gemini-Ultra \citep{gemini}, remain inaccessible to the public, while the available open-source alternatives demonstrate significantly inferior performance.

This paper presents DeepSeekMath, a specialized language model that surpasses existing open-source models in mathematical reasoning and closely rivals GPT-4 on established academic benchmarks. Central to this advancement is the development of the DeepSeekMath Corpus, a vast and meticulously curated pre-training dataset containing 120B math tokens. The corpus is compiled from Common Crawl (CC) using a classifier based on fastText \citep{joulin2016fasttext}. For its initial training, the classifier is supplied with positive samples from OpenWebMath \citep{openwebmath}, while a varied assortment of unrelated web documents are designated as negative examples. After this first pass, the classifier is utilized to extract new potential math-related instances from CC, which are then validated and improved through manual annotation. The enhanced data are incorporated to retrain the classifier, leading to further improvements in accuracy. Our assessments confirm the corpus’s exceptional quality: the DeepSeekMath-Base 7B model attains 64.2\% accuracy on GSM8K \citep{gsm8k} and 36.2\% on the advanced-level MATH dataset \citep{MATH}, exceeding the performance of Minerva 540B \citep{minerva}. Additionally, the multilingual aspect of the DeepSeekMath Corpus contributes to better outcomes on Chinese mathematics benchmarks \citep{wei2023cmath,agieval}. We assert that our strategies for mathematical dataset creation constitute a valuable foundation for future inquiry, and anticipate further progress in subsequent work.

DeepSeekMath-Base is built upon DeepSeek-Coder-Base-v1.5 7B \citep{deepseek-coder}, given our findings that initializing from a code-centric pretrained model yields superior outcomes compared to using a standard large language model. Additionally, our experiments reveal that mathematical training augments performance on the MMLU \citep{mmlu} and BBH \citep{bbh} benchmarks. This suggests that such training not only strengthens mathematical problem-solving but also broadens the model’s general reasoning proficiency.

Following the pre-training phase, we perform mathematical instruction tuning on DeepSeekMath-Base utilizing datasets featuring chain-of-thought \citep{cot}, program-of-thought \citep{pot,pal}, and tool-integrated reasoning \citep{tora} approaches. The fine-tuned model, DeepSeekMath-Instruct 7B, surpasses all other 7B models and demonstrates performance on par with 70B instruction-tuned open-source counterparts.

In addition, we propose Group Relative Policy Optimization (GRPO), an alternative RL algorithm to Proximal Policy Optimization (PPO) \citep{schulman2017proximal}. Unlike PPO, GRPO omits the critic network, instead utilizing group-based score aggregation as the baseline, which leads to a notable reduction in computational overhead during training. Relying on a limited portion of English instruction tuning datasets, GRPO delivers marked improvements over the robust DeepSeekMath-Instruct model across both in-domain benchmarks (GSM8K: 82.9\% $\rightarrow$ 88.2\%, MATH: 46.8\% $\rightarrow$ 51.7\%) and on out-of-domain evaluations (for instance, CMATH: 84.6\% $\rightarrow$ 88.8\%) during RL fine-tuning. Additionally, we establish a unified framework for interpreting various related approaches—such as Rejection Sampling Fine-Tuning (RFT) \citep{yuan2023scaling}, Direct Preference Optimization (DPO) \citep{dpo}, PPO, and GRPO—highlighting that these strategies can be seen as manifestations of either direct or simplified forms of reinforcement learning. Our study involves comprehensive empirical analyses, including comparisons of online versus offline training, output-based versus process-oriented supervision, and investigations into both single-step and iterative RL, to thoroughly examine critical factors underpinning this unified view. Finally, we provide an explanation for the observed gains in instruction-tuned model performance via RL, and outline future research avenues aimed at enhancing RL efficacy within this cohesive paradigm.

\subsection{Contributions}
We present a scalable approach to math pre-training, and further investigate and analyze the application of reinforcement learning.

\textbf{Math Pre-Training at Scale}

\begin{itemize}[topsep=0pt]
    \item In our study, we demonstrate that the openly available Common Crawl dataset holds significant potential for use in mathematical applications. Through the application of a carefully crafted data selection process, we curate the DeepSeekMath Corpus—a robust dataset comprising 120B tokens extracted from web pages specifically filtered for mathematical relevance. This corpus substantially surpasses previous efforts, amounting to nearly 7 times the volume of the math web pages utilized by Minerva \citep{minerva} and 9 times greater than the recently introduced OpenWebMath \citep{openwebmath}.

\item DeepSeekMath-Base 7B, our pre-trained base model, demonstrates performance on par with Minerva 540B \citep{minerva}. This result suggests that model size alone does not solely determine proficiency in mathematical reasoning, and that robust results can be obtained by training a comparatively smaller model with high-quality data.

\item We present the results of our math training studies. Conducting code training before math training enhances the models' performance on mathematical tasks, regardless of whether tools are involved. This provides some insight into the enduring question: \textit{does code training improve reasoning abilities?} Our evidence suggests that it does, specifically in the context of mathematical reasoning.

\item While utilizing arXiv papers for training is a standard practice, particularly within the mathematical research community, our findings indicate that it does not lead to measurable gains across any of the mathematical benchmarks evaluated in this study.
\end{itemize}

\textbf{Exploration and Analysis of Reinforcement Learning}

\begin{itemize}[topsep=0pt]
    \item This work presents Group Relative Policy Optimization (GRPO), a reinforcement learning method characterized by both its efficacy and computational efficiency. Unlike conventional approaches, GRPO dispenses with the use of a critic model, opting instead to estimate baselines via aggregated group scores, thereby yielding a substantial decrease in training overhead relative to Proximal Policy Optimization (PPO).
    \item Empirical evidence shows that applying GRPO markedly improves the instruction-tuning performance of our model, DeepSeekMath-Instruct, leveraging only instruction-tuning datasets. Additionally, reinforcement learning via GRPO leads to notable gains in out-of-domain generalization.
    \item We establish a comprehensive framework for interpreting a variety of methods—including RFT, DPO, PPO, and GRPO—within a shared context. To further elucidate the core factors of this unified framework, we perform thorough experimental studies on facets such as online versus offline training, comparison between outcome and process-based supervision, and the efficacy of single-turn versus iterative reinforcement learning.
    \item Building on insights from our unified paradigm, we analyze the underlying contributors to the success of reinforcement learning for LLMs and outline several promising avenues for achieving further improvements in reinforcement learning performance.
\end{itemize}

\subsection{Summary of Evaluations and Metrics}

\begin{itemize}[topsep=0pt]
    \item \textbf{English and Chinese Mathematical Reasoning}: Our evaluation rigorously measures model performance on both English and Chinese mathematical reasoning tasks, spanning problem sets from elementary to collegiate complexity. For English, the benchmarks consist of GSM8K \citep{gsm8k}, MATH \citep{MATH}, SAT \citep{llemma}, OCW Courses \citep{minerva}, and MMLU-STEM \citep{mmlu}. The Chinese evaluation utilizes MGSM-zh \citep{mgsm}, CMATH \citep{wei2023cmath}, Gaokao-MathCloze \citep{agieval}, and Gaokao-MathQA \citep{agieval}. Assessment includes both the generation of standalone, fully detailed textual solutions (without recourse to external tools), as well as solving tasks via Python programming.

On English evaluation tasks, DeepSeekMath-Base demonstrates competitive performance with the proprietary Minerva 540B \citep{minerva}, and consistently outperforms all open-source base models, such as Mistral 7B \citep{mistral} and Llemma-34B \citep{llemma}, irrespective of their math pre-training background—frequently by a considerable margin. Furthermore, DeepSeekMath-Base achieves higher results on Chinese benchmarks, which can be attributed to our decision not to limit the math pre-training corpus to English data, as was done in earlier research \citep{minerva,llemma}, but rather to incorporate a broad selection of high-quality, non-English content. Through the integration of mathematical instruction tuning and reinforcement learning, DeepSeekMath-Instruct and DeepSeekMath-RL exhibit notable advances, with both surpassing the 50\% accuracy threshold on the challenging, competition-level MATH dataset—a first for open-source models.

\item \textbf{Formal Mathematics}: DeepSeekMath-Base is assessed on its ability to autoformalize informal mathematical statements using the informal-to-formal theorem proving benchmark from \citep{dsp_proof}, specifically on the miniF2F dataset \citep{minif2f} with Isabelle \citep{isabelle} as the selected proof assistant. The model achieves notable results in few-shot autoformalization scenarios.
\item \textbf{Natural Language Understanding, Reasoning, and Code}: To holistically measure the model’s competence in general language comprehension, logical reasoning, and programming, DeepSeekMath-Base is evaluated on the Massive Multitask Language Understanding (MMLU) benchmark \citep{mmlu}, spanning 57 multiple-choice tasks across a wide array of disciplines. Additionally, its reasoning abilities are examined using BIG-Bench Hard (BBH) \citep{bbh}, which includes 23 complex, multi-step reasoning problems, and its coding proficiency is gauged through HumanEval \citep{codex} and MBPP \citep{mbpp}, both standard benchmarks for code generation tasks. Math pre-training confers observable improvements in both understanding and reasoning abilities.

\section{Math Pre-Training}

\subsection{Data Collection and Decontamination} This section describes the methodology employed to build the DeepSeekMath Corpus utilizing Common Crawl data. As illustrated in Figure \ref{fig:data_collect}, an iterative pipeline is employed, showcasing a stepwise method for extracting a comprehensive mathematical dataset from Common Crawl. The process initiates from a seed corpus, which consists of a limited yet high-quality set of math-focused data. It should be emphasized that this strategy is generalizable and can be adapted for other areas, including programming.

Initially, we select OpenWebMath \citep{openwebmath}—a curated dataset of high-quality mathematical resources from the web—as our foundational seed corpus. Leveraging this dataset, we train a fastText model \citep{joulin2016fasttext} to identify additional web pages resembling those in OpenWebMath. For model training, we draw a random sample of 500,000 entries from the seed corpus to serve as positive examples, and another 500,000 pages from Common Crawl to act as negative samples. Training is conducted using an open-source fastText library\footnote{\url{https://fasttext.cc}}, with the model configured to a vector dimensionality of 256, a learning rate set to 0.1, a maximum word n-gram length of 3, a minimum threshold of 3 word occurrences, and 3 epochs for training. To compress the Common Crawl dataset, we apply URL deduplication and near-duplicate detection methods, reducing the dataset to 40 billion HTML web pages. The trained fastText model is then applied to the deduplicated Common Crawl for retrieving candidate mathematical pages. To eliminate low-quality mathematics content, we rank the retrieved pages according to the fastText-generated scores and retain only the highest scoring subset. The final data size is further validated through a series of pre-training experiments using the top 40B, 80B, 120B, and 160B tokens. For the initial phase, we opt to retain the leading 40B tokens.

Following the initial round of data harvesting, a significant number of mathematical web pages remain uncollected, primarily due to limited diversity within the positive examples used to train the fastText model. To address this limitation, we expand the seed corpus by sourcing additional mathematical web domains, thereby facilitating more robust optimization of the fastText model. Concretely, we first partition the entire Common Crawl dataset by domain, where each domain comprises all web pages with a shared base URL. For every domain, we determine the fraction of web pages captured during the first data collection pass. Domains with more than 10\% of their content gathered are deemed likely to contain mathematics-related material (for instance, \url{mathoverflow.net}). 
We then perform manual annotation to pinpoint URLs within these domains that are indicative of mathematical content (such as \url{mathoverflow.net/questions}). Any yet uncollected web pages corresponding to these URLs are incorporated into the seed corpus. This methodology allows for the extraction of additional positive instances, which in turn facilitates the retraining of a more effective fastText model, capable of recalling a greater volume of mathematical data in the next data collection cycle. Upon completing four rounds of data collection, we accumulate 35.5M mathematical web pages, summing to 120B tokens. During the fourth round, it becomes apparent that about 98\% of pages were already retrieved in the previous iteration, leading us to conclude the data collection process.

In order to prevent contamination of our benchmarks, we adopt the approach outlined in \cite{deepseek-coder}, excluding web pages that feature either questions or answers sourced from prominent English mathematical benchmarks—including GSM8K \citep{gsm8k} and MATH \citep{MATH}—as well as Chinese benchmarks such as CMATH \citep{wei2023cmath} and AGIEval \citep{agieval}. The filtering process is defined as follows: if a segment of text contains any exact 10-gram overlap with substrings from the aforementioned evaluation benchmarks, that segment is eliminated from our mathematics training dataset. In cases where benchmark text comprises fewer than 10 grams but consists of at least 3 grams, we perform an exact-match filter to remove affected web pages and mitigate contamination.

\subsection{Validating the Quality of the DeepSeekMath~Corpus}

We run pre-training experiments to investigate how the DeepSeekMath~Corpus is compared with the recently released math-training corpora:

\begin{itemize}[topsep=0pt]
    \item \textbf{MathPile} \citep{mathpile}: This dataset is comprised of 8.9 billion tokens from various sources such as textbooks, Wikipedia, ProofWiki, CommonCrawl, StackExchange, and arXiv, with arXiv providing more than 85\% of the total content.
    \item \textbf{OpenWebMath} \citep{openwebmath}: This collection extracts mathematical material from CommonCrawl, amounting to 13.6 billion tokens after filtering for mathematical relevance.
    \item \textbf{Proof-Pile-2} \citep{llemma}: Incorporating OpenWebMath, 10.3 billion tokens from AlgebraicStack (focused on mathematical programming), and 28.0 billion tokens from arXiv articles, this mathematical text repository is used according to \cite{llemma} with a token ratio between arXiv, web, and code of 2:4:1 during our evaluations.
\end{itemize}

\subsubsection{Training Setting}
\label{sec:quality-policy}
We fine-tune a foundational language model with 1.3B parameters—built upon the DeepSeek LLMs architecture \citep{deepseek-llm} and referred to as DeepSeek-LLM 1.3B—specifically for mathematics by training it on each mathematical dataset independently for 150B tokens. All model training is performed using the streamlined and efficient HAI-LLM \citep{haillm} framework. In accordance with DeepSeek LLMs' optimization strategies, we employ the AdamW optimizer \citep{adamW}, configuring it with $\beta_1=0.9$, $\beta_2=0.95$, and $\mathrm{weight\_decay}=0.1$. The learning rate follows a stepped schedule: after 2,000 warmup steps, it climbs to its maximum, then drops to 31.6\% of its peak at 80\% completion of training, and subsequently tapers further to 10.0\% of the peak at the 90\% training mark. The learning rate is capped at a maximum of 5.3e-4, and training utilizes a batch size of 4M tokens with sequences up to 4K tokens in length.

\subsubsection{Evaluation Results}

\textbf{The DeepSeekMath~Corpus is of high quality, covers multilingual mathematical content, and is the largest in size.}

\begin{itemize}[topsep=0pt]
    \item \textbf{High-quality}:
    We assess model capabilities across 8 mathematical benchmarks utilizing few-shot chain-of-thought prompting \cite{cot}. As indicated in Table \ref{tab:corpora_comparison}, the model trained with the DeepSeekMath~Corpus consistently outperforms competitors. Furthermore, Figure \ref{fig:corpora_comparisons} illustrates that, at the 50B token mark (corresponding to a single complete pass over Proof-Pile-2), the DeepSeekMath-trained model surpasses its Proof-Pile-2 counterpart, reflecting the superior average quality inherent to the DeepSeekMath Corpus.
    \item \textbf{Multilingual}:
    The DeepSeekMath~Corpus provides a diverse linguistic dataset, with English and Chinese constituting the primary languages. According to Table \ref{tab:corpora_comparison}, training with the DeepSeekMath~Corpus results in increased mathematical reasoning proficiency for both English and Chinese tasks. Conversely, the available mathematical datasets, which are chiefly composed in English, contribute little and can even negatively impact performance on Chinese-language mathematical reasoning.
    \item \textbf{Large-scale}:
    In terms of size, the DeepSeekMath~Corpus far exceeds current mathematical datasets. Figure \ref{fig:corpora_comparisons} shows that training DeepSeek-LLM 1.3B with the DeepSeekMath~Corpus produces a sharper learning trajectory and more persistent gains. By comparison, the baseline corpora are notably smaller and subject to repeated cycles through the data during training, causing their associated model improvements to plateau more rapidly.
\end{itemize}

\subsection{Training and Evaluating DeepSeekMath-Base 7B}

This section details DeepSeekMath-Base 7B, a foundational model characterized by advanced mathematical reasoning capabilities. The model builds upon DeepSeek-Coder-Base-v1.5 7B \citep{deepseek-coder} as its initialization and undergoes training over 500B tokens. The training corpus is composed of 56\% DeepSeekMath Corpus, 4\% AlgebraicStack, 10\% arXiv documents, 20\% Github code, with the final 10\% consisting of Common Crawl natural language data in both English and Chinese. The primary training configuration mirrors the approach discussed in Section \ref{sec:quality-policy}, with two notable exceptions: the learning rate peaks at 4.2e-4 and the batch size is increased to 10M tokens.

We undertake an in-depth evaluation of the mathematical proficiency of DeepSeekMath-Base 7B, emphasizing its competence in independently generating complete mathematical solutions, its effectiveness when utilizing external tools to address mathematical tasks, and its capacity for formal theorem verification. In addition to its mathematical evaluation, we present a broader characterization of the base model by examining its aptitude in natural language comprehension, logical reasoning, and programming abilities.

\paragraph{Mathematical Problem Solving with Step-by-Step Reasoning}

We assess the capabilities of DeepSeekMath-Base in addressing mathematical questions via few-shot chain-of-thought prompting \citep{cot}, utilizing eight benchmarks presented in both English and Chinese. These benchmarks feature a range of quantitative reasoning tasks—such as GSM8K \citep{gsm8k}, MATH \citep{MATH}, and CMATH \citep{wei2023cmath}—as well as multiple-choice formats, exemplified by MMLU-STEM \citep{mmlu} and Gaokao-MathQA \citep{agieval}. The benchmarks collectively span various mathematical domains and encompass problem difficulties from elementary through collegiate levels.

Table \ref{tab:base_math_cot} highlights that DeepSeekMath-Base 7B achieves the highest scores across all eight benchmarks when compared to other open-source base models, including the popular Mistral 7B \citep{mistral} and the newer Llemma 34B \citep{llemma}, which has been trained on the Proof-Pile-2 dataset \citep{llemma}. Remarkably, DeepSeekMath-Base exceeds other open-source base models by more than 10\% absolute margin on the MATH dataset, a benchmark representing competition-level problems, and even outperforms Minerva 540B \citep{minerva}—a proprietary model 77 times its size built upon PaLM \citep{palm} and additionally trained on mathematical corpora.

\paragraph{Mathematical Problem Solving with Tool Use} We assess the effectiveness of leveraging program-assisted mathematical reasoning on both GSM8K and MATH benchmarks via few-shot program-of-thought prompting techniques \citep{pot,pal}. For each question, models are instructed to generate a Python script to arrive at the solution, with access to libraries like \textit{math} and \textit{sympy} to support complex calculations. The output of the executed script is treated as the solution to the problem. According to the results summarized in Table \ref{tab:base_math_pot_proof}, DeepSeekMath-Base 7B surpasses the former state-of-the-art model, Llemma 34B.

\paragraph{Formal Mathematics}

The automation of formal proofs plays a significant role in promoting the correctness and dependability of mathematical arguments while also improving productivity, a trend that has garnered substantial scholarly interest in recent years. In this work, we assess the capabilities of DeepSeekMath-Base 7B on the informal-to-formal proving task introduced by \citep{dsp_proof}, which involves producing a formal proof given an informal statement, its formal version, and an informal proof. Our experiments utilize miniF2F \citep{minif2f}, a prominent benchmark focused on formalizing Olympiad-level mathematics, where we employ few-shot prompting to construct formal proofs in Isabelle for each test item. In alignment with the methodology of \cite{dsp_proof}, our approach instructs the model to output proof sketches, subsequently employing the standard automated tool Sledgehammer \citep{sledgehammer} to address any omitted justifications. The results, presented in Table \ref{tab:base_math_pot_proof}, indicate that DeepSeekMath-Base 7B achieves notable effectiveness in the task of proof autoformalization.

\paragraph{Natural Language Understanding, Reasoning, and Code}

We assess the natural language understanding capabilities of our model using the MMLU benchmark \citep{mmlu}, test reasoning abilities with BBH \citep{bbh}, and evaluate coding proficiency on both HumanEval \citep{codex} and MBPP \citep{mbpp}. Table \ref{tab:understanding_reasoning_code} demonstrates that DeepSeekMath-Base 7B achieves notable improvements in MMLU and BBH scores when compared to its predecessor, DeepSeek-Coder-Base-v1.5 \citep{deepseek-coder}, thereby confirming that targeted math-oriented training benefits both understanding and reasoning. Furthermore, the addition of code tokens during ongoing training enables DeepSeekMath-Base 7B to preserve the coding performance observed in DeepSeek-Coder-Base-v1.5 across the coding evaluations. In summary, DeepSeekMath-Base 7B consistently and substantially surpasses the general-purpose Mistral 7B model \citep{mistral} on all three benchmarks assessing reasoning and coding.

\section{Supervised Fine-Tuning}

\subsection{SFT Data Curation}

We develop a comprehensive instruction-tuning dataset for mathematics that encompasses both English and Chinese problem statements. The dataset spans a variety of mathematical disciplines and includes tasks across multiple difficulty tiers. Each problem is matched with a solution articulated in formats such as chain-of-thought (CoT) \citep{cot}, program-of-thought (PoT) \citep{pot,pal}, and tool-integrated reasoning \citep{tora}. Altogether, the dataset comprises 776K training instances.

\begin{itemize}[topsep=0pt]
    \item \textbf{English mathematical datasets}: Our work provides tool-augmented solution annotations for both GSM8K and MATH datasets. Additionally, we utilize a curated sample of MathInstruct \citep{MathInstruct} and leverage the training portion of Lila-OOD \citep{lila}, which consists of problems answered using CoT or PoT methodologies. The English datasets span a wide array of mathematical domains, including algebra, probability, number theory, calculus, and geometry.
    \item \textbf{Chinese mathematical datasets}: We assemble a collection of K-12 Chinese mathematical problems, representing 76 distinct sub-domains such as linear equations, each paired with annotations using both CoT reasoning and tool-integrated approaches.
\end{itemize}

\subsection{Training and Evaluating DeepSeekMath-Instruct 7B}

This section presents DeepSeekMath-Instruct 7B, which is fine-tuned for mathematical instructions using DeepSeekMath-Base as its foundation. Training samples are combined in a random order up to a maximum context window of 4,000 tokens. The model undergoes 500 training steps, utilizing a batch size of 256 and a fixed learning rate of 5e-5.

We evaluate models' mathematical performance both without and with tool use, on 4 quantitative reasoning benchmarks in English and Chinese. We benchmark our model against the leading models of the time:

\begin{itemize}[topsep=0pt]
    \item \textbf{Closed-source models} encompass: (1) the GPT series, with GPT-4 \citep{gpt4} and GPT-4 Code Interpreter~\footnote{\url{https://openai.com/blog/chatgpt-plugins\#\#code-interpreter}} recognized as the most advanced, (2) Gemini Ultra and Pro \citep{gemini}, (3) Inflection-2 \citep{inflection-2}, (4) Grok-1~\footnote{\url{https://x.ai/model-card}},
    and several newer offerings from Chinese organizations, such as (5) Baichuan-3~\footnote{\url{https://www.baichuan-ai.com}},
    and (6) the most recent GLM-4~\footnote{\url{https://open.bigmodel.cn/dev/api\#glm-4}}

    from the GLM family \citep{glm}.

These models serve broad applications, and the majority have been refined through various alignment strategies.

\item \textbf{Open-source models} encompass: general-purpose systems such as (1) DeepSeek-LLM-Chat 67B \citep{deepseek-llm}, (2) Qwen 72B \citep{qwen}, (3) SeaLLM-v2 7B \citep{seallm}, and (4) ChatGLM3 6B \citep{chatglm3}. Additionally, models optimized for mathematics include: (5) InternLM2-Math 20B~\footnote{\url{https://github.com/InternLM/InternLM-Math}}, which is derived from InternLM2 and further trained on mathematical data followed by instruction tuning, (6) Math-Shepherd-Mistral 7B, utilizing PPO training \citep{schulman2017proximal} with a process-supervised reward signal applied to Mistral 7B \citep{mistral}, (7) the WizardMath collection \citep{wizardmath}, designed to enhance mathematical capabilities in Mistral 7B and Llama-2 70B \citep{llama2} via evolve-instruct—a variant of instruction tuning with AI-generated instructions—and PPO optimization, leveraging datasets mainly from GSM8K and MATH, (8) MetaMath 70B \citep{metamath}, which consists of Llama-2 70B that has been fine-tuned on an expanded version of GSM8K and MATH, (9) ToRA 34B \cite{tora}, featuring CodeLlama 34B adapted for tool-based mathematical reasoning, and (10) MAmmoTH 70B \citep{MathInstruct}, in which Llama-2 70B is tuned with instructions from MathInstruct.

Referencing Table \ref{tab:sft_rl_math}, in the evaluation scenario where tool use is not permitted, DeepSeekMath-Instruct 7B exhibits robust capabilities in step-by-step reasoning tasks. Remarkably, on the competition-level MATH dataset, our model exceeds the performance of every open-source alternative as well as most commercial offerings (such as Inflection-2 and Gemini Pro) by a margin of at least 9\% in absolute terms. This advantage persists even over much larger models (e.g., Qwen 72B) or those that have undergone additional training with math-centric reinforcement learning (like WizardMath-v1.1 7B). DeepSeekMath-Instruct performs on par with leading Chinese proprietary models, GLM-4 and Baichuan-3, on the MATH benchmark. Nonetheless, it does not yet match the results achieved by GPT-4 and Gemini Ultra.

When assessed under conditions permitting both natural language reasoning and the utilization of programmatic tools for problem solving, DeepSeekMath-Instruct 7B attains nearly 60\% accuracy on the MATH benchmark, outperforming all other open-source models currently available. Across additional benchmarks, our model achieves results comparable to those of DeepSeek-LLM-Chat 67B, the previous leading approach, despite being only one-tenth its size.

\subsection{Group Relative Policy Optimization} Following the Supervised Fine-Tuning (SFT) phase, reinforcement learning (RL) has demonstrated significant utility in enhancing large language models' capacity for mathematical reasoning \citep{wang2023math,wizardmath}. In what follows, we present our novel RL approach, Group Relative Policy Optimization (GRPO), designed to offer both efficiency and effectiveness.

\subsubsection{From PPO to GRPO} Proximal Policy Optimization (PPO) \citep{schulman2017proximal} is an actor-critic RL algorithm that is widely used in the RL fine-tuning stage of LLMs \citep{ouyang2022training}. In particular, it optimizes LLMs by maximizing the following surrogate objective:

\begin{equation}
\footnotesize
    \mathcal{J}_{PPO}(\theta) = \mathbb{E}{[q \sim P(Q),\ o \sim \pi_{\theta_{old}}(O|q)]} \frac{1}{|o|} \sum_{t=1}^{|o|} \min \left[ \frac{\pi_\theta(o_{t} | q, o_{<t})}{\pi_{\theta_{old}}(o_{t} | q, o_{<t})} A_{t},\ \text{clip} \left( \frac{\pi_\theta(o_{t} | q, o_{<t})}{\pi_{\theta_{old}}(o_{t} | q, o_{<t})},\ 1 - \varepsilon,\ 1 + \varepsilon \right) A_{t} \right],
\end{equation}

Here, $\pi_{\theta}$ denotes the current policy model, while $\pi_{\theta_{old}}$ refers to the previous policy model. The symbols $q$ and $o$ represent questions and responses, respectively, which are sampled from the dataset of questions and generated by the old policy $\pi_{\theta_{old}}$. The hyperparameter $\varepsilon$ is introduced in PPO as a mechanism to stabilize the learning process via clipping. The term $A_t$ indicates the advantage at time step $t$, which is estimated through Generalized Advantage Estimation (GAE) \citep{gae} utilizing the sequence of rewards $\{r_{\ge t}\}$ along with the estimated value function $V_{\psi}$. Consequently, PPO requires the concurrent training of both the policy and value function. To prevent excessive reward model optimization, it is customary to include a per-token KL divergence penalty between the output distribution of the policy and a fixed reference model in the reward at each token \citep{ouyang2022training}, that is,

\begin{equation}
    r_{t} = r_\varphi(q, o_{\le t}) - \beta \log\frac{\pi_{\theta}(o_{t}|q, o_{<t})}{\pi_{ref}(o_{t}|q, o_{<t})},
\label{eq:PPO-reward}
\end{equation}

Here, $r_\varphi$ denotes the reward model, while $\pi_{ref}$ signifies the reference model—most often represented by the initial SFT model. The term $\beta$ serves as the scaling factor for the KL divergence penalty.

As the value function employed in PPO is typically another model of comparable size as the policy model, it brings a substantial memory and computational burden. Additionally, during RL training, the value function is treated as a baseline in the calculation of the advantage for variance reduction. While in the LLM context, usually only the last token is assigned a reward score by the reward model, which may complicate the training of a value function that is accurate at each token. To address this, as shown in Figure \ref{fig:grpo}, we propose Group Relative Policy Optimization (GRPO), which obviates the need for additional value function approximation as in PPO, and instead uses the average reward of multiple sampled outputs, produced in response to the same question, as the baseline. More specifically, for each question $q$, GRPO samples a group of outputs $\{o_1, o_2, \cdots, o_G\}$  from the old policy  $\pi_{\theta_{old}}$  and then optimizes the policy model by maximizing the following objective:

\begin{equation}
\footnotesize
\begin{split}
    \mathcal{J}_{GRPO}(\theta) &= \mathbb{E}{[q \sim P(Q), \{o_i\}_{i=1}^G \sim \pi_{\theta_{old}}(O|q)]}  \\
    & \frac{1}{G}\sum_{i=1}^G\frac{1}{|o_i|} \sum_{t=1}^{|o_i|} \left\{ \min \left[ \frac{\pi_\theta(o_{i,t} | q, o_{i,<t})}{\pi_{\theta_{old}}(o_{i,t} | q, o_{i,<t})} \hat{A}_{i,t}, \text{clip} \left( \frac{\pi_\theta(o_{i,t} | q, o_{i,<t})}{\pi_{\theta_{old}}(o_{i,t} | q, o_{i,<t})}, 1 - \varepsilon, 1 + \varepsilon \right)  \hat{A}_{i,t} \right] - \beta \mathbb{D}_{KL}\left[\pi_{\theta} || \pi_{ref}\right]\right\} ,
\end{split}
\label{eq:GRPO-obj}
\end{equation}

This objective function for $\mathcal{J}_{GRPO}(\theta)$ consists of an expectation over questions $q$ sampled from $P(Q)$ and $G$ generations $\{o_i\}$ drawn from the previous policy $\pi_{\theta_{old}}(O|q)$. For each generation $o_i$, it averages across all timesteps $t$ within the sequence. The core expression within the curly braces takes the minimum of two terms: the advantage-weighted probability ratio between the current and old policy, and a clipped version of this ratio (bounded by $1 \pm \varepsilon$), each scaled by the estimated advantage $\hat{A}_{i,t}$. Additionally, there is a regularization component that penalizes the Kullback-Leibler divergence between the updated policy $\pi_\theta$ and a reference policy $\pi_{ref}$, scaled by the hyperparameter $\beta$.

where $\varepsilon$ and $\beta$ are hyper-parameters, and $\hat{A}_{i,t}$  is the advantage calculated based on relative rewards of the outputs inside each group only, which will be detailed in the following subsections. The group relative way that GRPO leverages to calculate the advantages, aligns well with the comparative nature of rewards models,  as reward models are typically trained on datasets of comparisons between outputs on the same question. Also note that, instead of adding KL penalty in the reward, GRPO regularizes by directly adding the KL divergence between the trained policy and the reference policy to the loss, avoiding complicating the calculation of  $\hat{A}_{i,t}$. And different from the KL penalty term used in (\ref{eq:PPO-reward}), we estimate the KL divergence with the following unbiased estimator \citep{kl_approx}:

\begin{equation}
\small
    \mathbb{D}_{KL}\left[\pi_{\theta} || \pi_{ref}\right] = \frac{\pi_{ref}(o_{i,t}|q,o_{i,<t})}{\pi_{\theta}(o_{i,t}|q,o_{i,<t})}- \log\frac{\pi_{ref}(o_{i,t}|q,o_{i,<t})}{\pi_{\theta}(o_{i,t}|q,o_{i,<t})} - 1,
\end{equation}

which is guaranteed to be positive.

\begin{algorithm}[t]
  \small
  \caption{Iterative Group Relative Policy Optimization}
  \textbf{Input} initial policy model $\pi_{\theta_{\text{init}}}$; reward models $r_\varphi$; task prompts $\mathcal{D}$; 
  hyperparameters $\varepsilon$, $\beta$, $\mu$
  \begin{algorithmic}[1]
    \State Set policy model $\pi_\theta$ to $\pi_{\theta_{\text{init}}}$
    \For{iteration = 1, \dots, I}
      \State Assign current reference model $\pi_{ref} \gets \pi_{\theta}$
      \For{step = 1, \dots, M}
        \State Draw a minibatch $\mathcal{D}_b$ from $\mathcal{D}$
        \State Record current policy $\pi_{\theta_{old}} \gets \pi_{\theta}$
        \State For each prompt $q$ in $\mathcal{D}_b$, sample $G$ outputs $\{o_i\}_{i=1}^G \sim \pi_{\theta_{old}} (\cdot \mid q) $
        \State For every output $o_i$ in the set, compute corresponding rewards $\{r_i\}_{i=1}^{G}$ utilizing $r_{\varphi}$
        \State Calculate $\hat{A}_{i,t}$, the group relative advantage for the $t$-th token of every $o_i$
        \For{each GRPO update = 1, \dots, $\mu$}
          \State Update $\pi_{\theta}$ by optimizing the GRPO objective as given in Equation \ref{eq:GC-GRPO}
        \EndFor
      \EndFor 
      \State Enhance $r_\varphi$ with continued training leveraging a replay buffer. 
    \EndFor 
  \end{algorithmic}
  \textbf{Output} $\pi_\theta$
  \label{alg:iter-grpo}
\end{algorithm}

\subsubsection{Outcome Supervision RL with GRPO} Specifically, given a question $q$, a batch of outputs $\{o_1, o_2, \cdots, o_G\}$ is generated from the previous policy $\pi_{\theta_{old}}$. Each of these outputs is then evaluated using a reward model, producing a corresponding set of rewards $\mathbf{r}=\{r_1, r_2, \cdots, r_G\}$. To normalize these rewards, each reward is adjusted by removing the mean of the group and dividing by its standard deviation. In outcome supervision, this normalized reward is assigned at the conclusion of each output $o_i$, and the advantage for every token within the output, $\hat{A}_{i, t}$, is set equal to this normalized value: $\hat{A}_{i, t} = \widetilde{r}_i = \frac{r_i- {\rm mean}(\mathbf{r})}{{\rm std}(\mathbf{r})}$. The policy is then updated by maximizing the objective described in equation (\ref{eq:GRPO-obj}).

\subsubsection{Process Supervision RL with GRPO} While outcome supervision delivers a singular reward upon completion of each generated output, this approach may lack sufficiency and efficiency when guiding policies through intricate mathematical reasoning. In line with \cite{wang2023math}, we incorporate process supervision, wherein feedback is administered at each intermediate reasoning step rather than only at the conclusion. Specifically, for a given problem $q$ and $G$ sampled output sequences $\{o_1, o_2, \cdots, o_G\}$, a process reward model evaluates each reasoning step within the sequences, producing a set of rewards: $\mathbf{R} = \{ \{r_1^{index(1)},\cdots,r_1^{index(K_1)}\}, \cdots,  \{r_G^{index(1)},\cdots,r_G^{index(K_G)}\} \}$. Here, $index(j)$ designates the index corresponding to the terminal token of the $j$-th step, and $K_i$ denotes the total number of reasoning steps for the $i$-th sequence. We further perform normalization of these reward signals using the mean and standard deviation: $\widetilde{r}_i^{index(j)} = \frac{r_i^{index(j)} - {\rm mean(\mathbf{R})}}{{\rm std(\mathbf{R})}}$.

Next, process supervision determines the advantage for each token as the cumulative sum of normalized rewards assigned to subsequent steps, formally given by $\hat{A}_{i, t} = \sum_{index(j) \ge t} \widetilde{r}_i^{index(j)}$. Finally, this advantage estimation is utilized to refine the policy by maximizing the objective function specified in equation (\ref{eq:GRPO-obj}).

\subsubsection{Iterative RL with GRPO} As training in reinforcement learning advances, the earlier reward model may become inadequate for effectively guiding the increasingly capable policy model. To address this, we investigate an approach involving iterative RL using GRPO. Specifically, as outlined in Algorithm \ref{alg:iter-grpo}, the iterative GRPO procedure involves producing updated training datasets for the reward model through samples collected from the most recent policy model. The previous reward model is further refined using these datasets while a replay strategy retains 10\% of earlier samples. Subsequently, we update the reference model to correspond with the latest policy model and proceed to optimize the policy model using the updated reward model.

\subsection{Training and Evaluating DeepSeekMath-RL}

We utilize RL grounded on DeepSeekMath-Instruct 7B. The RL training utilizes chain-of-thought formatted problems drawn from the GSM8K and MATH datasets in the SFT corpus, encompassing approximately 144K items. To examine how RL influences benchmarks without extensive data exposure, we deliberately omit other SFT data during RL. The construction of the reward model training set mirrors the approach in \citep{wang2023math}. Our initial reward model is trained with DeepSeekMath-Base 7B, employing a learning rate of 2e-5. For GRPO, the policy model uses a learning rate of 1e-6, and the KL penalty coefficient is fixed at 0.04. We generate $64$ samples per problem. The maximum sequence length is capped at 1024 tokens, and we set the batch size to 1024 for training. The policy model undergoes a single update after each exploratory stage. Benchmark evaluations for DeepSeekMath-RL 7B are conducted using the same settings as for DeepSeekMath-Instruct 7B. For the DeepSeekMath-RL 7B model, tasks like GSM8K and MATH—both utilizing chain-of-thought reasoning—are categorized as in-domain, whereas all other benchmarks are classified as out-of-domain.

As shown in Table \ref{tab:sft_rl_math}, the results compare open-source and closed-source models employing both chain-of-thought and tool-integrated reasoning on English and Chinese datasets. The analysis reveals: 1) When applying chain-of-thought reasoning, DeepSeekMath-RL 7B achieves 88.2\% accuracy on GSM8K and 51.7\% on MATH. These scores not only exceed those of every open-source competitor within the 7B to 70B model class but also surpass most closed-source systems. 2) Notably, the DeepSeekMath-RL 7B model is exclusively trained on instruction tuning data in the chain-of-thought format from GSM8K and MATH, using DeepSeekMath-Instruct 7B as its foundation. Although its training data is comparatively limited, DeepSeekMath-RL 7B delivers consistently superior results over DeepSeekMath-Instruct 7B on every metric, underscoring the advantages introduced by reinforcement learning.

\section{Discussion}
This section presents our observations derived from both pre-training and reinforcement learning experiments.

\subsection{Lessons Learnt in Pre-Training}

We begin by discussing our observations during the pre-training phase. Unless explicitly stated, the experiments employ the configuration described in Section \ref{sec:quality-policy}. Additionally, throughout this section, the DeepSeekMath Corpus referenced corresponds to the 89B-token dataset compiled during the second round of data gathering.

\subsubsection{Code Training Benefits Mathematical Reasoning}

A widely circulated but unsubstantiated theory claims that engaging in code training enhances reasoning capabilities. In this work, we seek to provide evidence addressing this assertion specifically in mathematics: code training bolsters models’ proficiency in mathematical reasoning, regardless of whether they utilize auxiliary tools.

To study how code training affects mathematical reasoning, we experimented with the following two-stage training and one-stage training settings:

\textbf{Two-Stage Training}

\begin{itemize}[topsep=0pt]
    \item \textbf{Code Training for 400B Tokens $\rightarrow$ Math Training for 150B Tokens}: The DeepSeek-LLM 1.3B model undergoes pretraining on 400B code tokens, which is subsequently followed by an additional 150B tokens focusing on mathematics;
    \item \textbf{General Training for 400B Tokens $\rightarrow$ Math Training for 150B Tokens}: For comparison, we conduct a parallel experiment using 400B general tokens (sourced from a comprehensive general corpus assembled by DeepSeek-AI) in place of code tokens during the initial training phase, to explore whether code token pretraining brings particular benefits to the enhancement of mathematical reasoning.
\end{itemize}

\textbf{One-Stage Training}

\begin{itemize}[topsep=0pt]
    \item \textbf{Mathematics Pretraining with 150B Tokens}: DeepSeek-LLM 1.3B undergoes training on a dataset containing 150 billion mathematics-related tokens;
    \item \textbf{Combined Training Utilizing 400B Code Tokens and 150B Math Tokens}: Directly fine-tuning on math tokens after code training leads to diminished code capabilities. Therefore, we assess if concurrently training with both code and math tokens from the outset can enhance mathematical reasoning skills, while simultaneously mitigating the adverse effects of catastrophic forgetting on coding knowledge.
\end{itemize}

\paragraph{Results}
The downstream outcomes across various training configurations are presented in Table \ref{tab:code-math} and Table \ref{tab:code-math-general-eval}.

Incorporating code training proves advantageous for program-aided mathematical reasoning, whether employing a two-stage or one-stage training approach. According to Table \ref{tab:code-math}, in the two-stage training framework, code-focused training by itself leads to marked improvements in solving GSM8K and MATH tasks when utilizing Python. Introducing math training during the subsequent stage delivers additional gains in performance. Notably, with one-stage training, combining both code and math tokens in a single process helps to effectively address the problem of catastrophic forgetting observed with the two-stage strategy, and also reinforces both coding skills (Table \ref{tab:code-math-general-eval}) and the capacity for program-aided mathematical reasoning (Table \ref{tab:code-math}).

Training with code data likewise enhances mathematical reasoning even when no tools are used. In a two-stage training regimen, the first phase involving code yields notable, albeit moderate, improvements. Furthermore, it facilitates greater efficiency in later math-focused training, culminating in superior overall performance. In contrast, merging code and mathematical tokens during a single-stage training process appears to detract from math reasoning ability in the absence of tool usage. This outcome may be attributed to the constrained capacity of DeepSeek-LLM 1.3B, which may be insufficient to comprehensively absorb and integrate both code and mathematics data at the same time.

\subsubsection{ArXiv Papers Seem Ineffective in Improving Mathematical Reasoning}

ArXiv papers are commonly included as a component of math pre-training data \citep{minerva,gpt-f,llemma,mathpile}. However, detailed analysis regarding their impact on mathematical reasoning has not been extensively conducted. Perhaps counter-intuitively, according to our experiments, arXiv papers seem ineffective in improving mathematical reasoning. We experiment with models of different sizes, including DeepSeek-LLM 1.3B and DeepSeek-Coder-Base-v1.5 7B \citep{deepseek-coder}, using arXiv corpora that underwent varied processing pipelines:

\begin{itemize}[topsep=0pt]
    \item \textbf{MathPile} \citep{mathpile}: consists of 8.9 billion tokens assembled through heuristic methods for cleaning and filtering, with scientific arXiv papers accounting for over 85\% of the content;
    \item \textbf{ArXiv-RedPajama} \citep{redpajama}: comprises all arXiv LaTeX documents after removing preambles, comments, macros, and bibliography sections, resulting in a dataset of 28.0 billion tokens.
\end{itemize}

For our study, we conducted separate pre-training of DeepSeek-LLM 1.3B on 150B tokens and DeepSeek-Coder-Base-v1.5 7B on 40B tokens, each using arXiv-only corpora. The results indicate that utilizing arXiv papers alone does not contribute to gains in mathematical reasoning skills. Both language models, when exposed solely to arXiv data, either stagnate or exhibit declines in performance across a spectrum of mathematical evaluation tasks with varying degrees of difficulty. These tasks encompass quantitative reasoning datasets such as GSM8K and MATH (Table \ref{tab:arxiv-cot}), multiple-choice tests including MMLU-STEM (Table \ref{tab:arxiv-cot}), as well as formal mathematics assessments like miniF2F (Table \ref{tab:arxiv_atp}).

However, this conclusion has its limitations and should be taken with a grain of salt. We have not yet studied:

\begin{itemize}[topsep=0pt]
    \item How arXiv tokens influence specialized mathematical tasks beyond those evaluated here, for instance, the informalization of theorems, which involves translating formal statements or proofs into more informal language;
    \item The influence exerted by arXiv tokens when used in conjunction with alternative data sources;
    \item If the advantages conferred by incorporating arXiv papers become more apparent with models of greater scale.
\end{itemize}

Thus, further exploration is required, which we leave for future studies.

\subsection{Insights of Reinforcement Learning}
\subsubsection{Towards to a Unified Paradigm} This section introduces a cohesive framework for examining various training strategies, including SFT, RFT, DPO, PPO, and GRPO. We also present experiments aimed at investigating the components involved in this overarching paradigm. In general, the gradient of a training method with respect to the parameter $\theta$ can be expressed as:

\begin{equation}
    \nabla_{\theta}\mathcal{J}_{\textcolor{red}{\mathcal{A}}}(\theta) = \mathbb{E}[\underbrace{(q,o) \sim \textcolor{red}{\mathcal{D}}}_{Data \ Source}]\left( \frac{1}{|o|} \sum_{t=1}^{|o|}  \underbrace{GC_{{\mathcal{A}}}(q, o, t, \textcolor{red}{\pi_{{rf}}})}_{Gradient \ Coefficient}  \nabla_{\theta}\log \pi_{\theta}(o_t | q, o_{<t})\right).
\label{eq:objective}
\end{equation}

There exist three key components: 1) \textit{Data Source $\mathcal{D}$}, which determines the training data; 2) \textit{Reward Function $\pi_{{rf}}$}, which is the source of the training reward signal; 3) \textit{Algorithm $\mathcal{A}$}: which processes the training data and the reward signal to the gradient coefficient $GC$ that determines the magnitude of the penalty or reinforcement for the data. We analyze several representative methods based on such a unified paradigm:

\begin{itemize}
    \item  \textbf{Supervised Fine-tuning (SFT)}: SFT involves adapting a pretrained model by training it further on datasets curated and selected by humans for SFT purposes.
    \item \textbf{Rejection Sampling Fine-tuning (RFT)}: RFT continues the optimization of the SFT model by training on output responses, generated by the SFT model in response to SFT queries, that are subsequently filtered according to the accuracy of their answers.
    \item \textbf{Direct Preference Optimization (DPO)}: DPO improves upon the SFT model through additional training on enhanced outputs, produced by the SFT model, guided by a pairwise DPO loss function.
    \item \textbf{Online Rejection Sampling Fine-tuning (Online RFT)}: In contrast to standard RFT, Online RFT sets the initial policy model as the SFT model but repeatedly updates it with augmented outputs that are generated by the currently-updated, real-time policy model itself.
    \item \textbf{PPO/GRPO}: PPO/GRPO begins with the SFT model to initialize the policy, then applies reinforcement training using outputs generated by the evolving, real-time policy model.
\end{itemize}

Table \ref{tab:data_policy} provides an overview of the components involved in these methods. For an in-depth explanation of the derivation procedure, see Appendix \ref{app:analysis-rl}.

\paragraph{Observation about Data Source} The data sources are categorized into two distinct types: online sampling and offline sampling. Online sampling refers to scenarios where the training data is collected based on the outcomes generated by the current training policy model during exploration. In contrast, offline sampling utilizes data collected from the outputs of the initial SFT model. Methods such as RFT and DPO utilize offline sampling, whereas Online RFT and GRPO operate using the online sampling approach.

As illustrated in Figure \ref{fig:rl-analysis}, our results indicate that Online RFT delivers markedly better performance than RFT on both evaluated benchmarks. In the initial training phase, Online RFT and RFT display similar performance, but as training proceeds, Online RFT achieves a clear performance lead, underscoring the benefits of the online training approach. This observation aligns with expectations: at the outset, the actor's behavior closely mirrors that of the SFT model, resulting in only slight discrepancies in sampled data. Over time, as training progresses, the differences in data sampled from the actor become more pronounced, making real-time data sampling increasingly advantageous.

\paragraph{Observation about Gradient Coefficient} The algorithm utilizes the reward signal by translating it into the gradient coefficient, which is then used to refine the model parameters. In our study, we decompose the reward function into two distinct types: `Rule' and `Model.' The `Rule' variant involves evaluating responses solely on the basis of answer correctness, while the `Model' approach employs a trained reward model that assigns scores to responses, with the model itself being trained on rule-based judgments. Equations \ref{eq:GC-RFT} and \ref{eq:GC-GRPO} demonstrate a notable distinction between GRPO and Online RFT: GRPO adjusts its gradient coefficient dynamically according to the reward values produced by the reward model, enabling nuanced strengthening or weakening of responses according to their assessed merit. Conversely, Online RFT does not implement such a mechanism; it omits penalties for incorrect answers and applies consistent positive reinforcement to all responses deemed correct, regardless of differences in reward magnitude.

As illustrated in Figure \ref{fig:rl-analysis}, GRPO outperforms online RFT, underscoring the effectiveness of modifying both positive and negative gradient coefficients. Moreover, GRPO+PS achieves better results than GRPO+OS, which suggests that employing more precise, step-sensitive gradient coefficients is advantageous. Additionally, we investigate iterative RL by performing two iterative rounds in our experiments. Figure \ref{fig:iter-rl} reveals that applying iterative RL substantially enhances performance, with the most notable gains observed during the initial iteration.

\subsubsection{Why RL Works?} In this study, we apply reinforcement learning to a selected subset of instruction tuning data, resulting in notable improvements over the baseline instruction-tuned model. To elucidate the reasons behind RL’s effectiveness, we assess both Pass@K and Maj@K accuracy for the Instruct and RL-enhanced models across two evaluation benchmarks. As depicted in Figure \ref{fig:maj-pass}, the introduction of RL yields higher scores in Maj@K but does not alter Pass@K outcomes. This pattern suggests that reinforcement learning mainly strengthens the stability of the output distribution, implying that \textbf{the observed gains stem from promoting the selection of correct answers among the TopK outputs, rather than enhancing the model’s underlying abilities.} In line with this, \citep{wang2023making} have discussed a \textbf{misalignment problem} affecting reasoning tasks in SFT models, and have demonstrated that such models can see improvements in reasoning ability by employing a range of preference alignment approaches \citep{yuan2023rrhf,song2023preference,wang2023making}.

\subsubsection{How to Achieve More Effective RL?}
\label{sec:effec_rl}
Our results illustrate that RL performs effectively on mathematical reasoning problems. Furthermore, we introduce a unified framework that facilitates the understanding of various notable training approaches. In this framework, each method can be interpreted as an application of RL—either in its explicit or implicit forms. According to Equation \ref{eq:objective}, three fundamental elements are identified: Data Source, Algorithm, and Reward Function. We outline several prospective avenues for further research concerning these three components.

\paragraph{Data Source} The foundation of any training approach lies in its data source. For RL scenarios, this specifically denotes the unlabeled questions alongside their responses, which are produced by sampling from the policy model. In this study, our focus is restricted to utilizing questions derived from the instruction tuning phase, and we employ a straightforward nucleus sampling method to generate outputs. We hypothesize that this limitation may underlie the observation that our RL pipeline results in improvements primarily for Maj@K metrics. Looking forward, we intend to apply our RL pipeline to out-of-distribution question prompts, integrating \textbf{more sophisticated sampling (decoding) approaches}, such as tree-search-based techniques \citep{yao2023tree}. Moreover, the use of \textbf{efficient inference methods} \citep{xia-etal-2023-speculative,leviathan2023fast,kwon2023efficient,xia2024unlocking}, which significantly influence how effectively policy models can be explored, is also deemed critically important.

\paragraph{Algorithms} In reinforcement learning, algorithms utilize both data and reward signals to adjust the gradient coefficient, thus enabling updates to the model parameters. As per Equation \ref{eq:objective}, contemporary approaches generally \textbf{TRUST} the reward function's feedback when increasing or decreasing the conditional probability of particular tokens. Nevertheless, this reliance presupposes the reward signal’s dependability, which cannot always be guaranteed in practice, particularly in highly intricate tasks. For instance, despite meticulous annotation by expert annotators, the PRM800K datasets \citep{lightman2023let} still exhibit around a 20\% rate of annotation errors\footnote{\url{https://github.com/openai/prm800k/issues/12\#issuecomment-1728491852}}. Consequently, we propose to investigate reinforcement learning algorithms designed to be resilient in the face of noisy reward feedback. We anticipate that alignment strategies characterized as \textbf{WEAK-TO-STRONG} \citep{burns2023weak} could drive substantive innovation within learning algorithms.

\paragraph{Reward Function} The reward function serves as the primary training signal in RL, and it is typically implemented as a neural reward model. There are three central avenues for advancing reward models: 1) \textbf{Strategies to improve the generalization of the reward model.} A robust reward model should generalize well to unfamiliar prompts and novel decoding results; otherwise, RL may only reinforce the existing output distribution of LLMs rather than driving genuine improvements in their underlying abilities; 2) \textbf{Capturing and representing the reward model’s uncertainty.} Quantifying uncertainty has the potential to mediate between imperfect reward models and learning algorithms designed to leverage such weaknesses for progression; 3) \textbf{Developing process reward models that are both high-quality and resource-efficient,} enabling them to offer nuanced, step-wise feedback throughout reasoning tasks \citep{lightman2023let,wang2023math}.

\section{Conclusion, Limitation, and Future Work} This work introduces DeepSeekMath, a model that surpasses all publicly available systems on the MATH benchmark at the competition level, while coming close to the achievements of proprietary models. DeepSeekMath~builds upon the DeepSeek-Coder-v1.5 7B foundation and receives continual training over 500B tokens, of which 120B are math-related tokens extracted from Common Crawl, forming a substantial part of the corpus. Through comprehensive ablation analyses, we demonstrate that web-derived content presents a valuable source for high-quality mathematical data, whereas the utility of arXiv proved less significant than anticipated. We propose Group Relative Policy Optimization (GRPO), a novel adaptation of Proximal Policy Optimization (PPO), capable of markedly enhancing mathematical reasoning with reduced memory overhead. Our experimental findings confirm that GRPO remains effective, even for DeepSeekMath-Instruct 7B, which already achieves strong benchmark results. Finally, we offer a unified framework for interpreting related methods and highlight several promising avenues for future reinforcement learning research.

While DeepSeekMath~demonstrates strong results on quantitative reasoning tasks, its performance lags behind that of proprietary models in areas such as geometry and theorem proving. For example, our preliminary tests reveal that DeepSeekMath~struggles with questions involving triangles and ellipses, suggesting that the model's pre-training and fine-tuning phases might have included insufficient data from these domains. Moreover, due to limitations in model size, DeepSeekMath~is less adept at utilizing few-shot prompts compared to GPT-4; whereas GPT-4 benefits from few-shot examples and shows marked improvement, DeepSeekMath~exhibits comparable outcomes regardless of whether zero-shot or few-shot techniques are applied. Moving forward, we intend to enhance our curated data selection process to assemble a higher-quality pre-training dataset. Furthermore, we plan to investigate promising approaches (Section \ref{sec:effec_rl}) for more robust reinforcement learning applied to LLMs.

\end{document}